\fignature{etacar_diffim_slit.jpg}{0.97}
{$\eta$~Car light echoes. The left panel shows the positions of
$\eta$~Car and our images (white box), plotted on an image
in the light of 3 different emission lines:
oxygen (blue), hydrogen (green), and sulfur (red) (credit:
Nathan Smith, University of Arizona/NOAO/AURA/NSF).
%
%http://www.noao.edu/image_gallery/html/im0667.html
%This image shows a giant star-forming region in the southern sky known as the Carina Nebula 
%(NGC3372), combining the light from 3 different filters tracing emission from 
% oxygen (blue), hydrogen (green), and sulfur (red). The color is also representative of the 
% temperature in the ionized gas: blue is relatively hot and red is cooler. The Carina Nebula 
% is a good example of how very massive stars rip apart the molecular clouds that give birth 
% to them. The bright star near the center of the image is $\eta$~Car, which is one of the most 
% massive and luminous stars known. This picture is a composite of several exposures made with 
% the Curtis Schmidt telescope at the Cerro Tololo Interamerican Observatory. We also have a 
% broad-band optical image, approximately true color, made photographically with the Curtis Schmidt. 
%
The middle panels show the images obtained with the CTIO 4-m
Blanco telescope of a region $\sim$0.5~degrees to the south of
$\eta$~Car at epochs 2003 March 10 ({\it A}), 2010 May 10 ({\it B}),
and 2011 Feb 6 ({\it C}), from top to bottom.  The right panels show
the difference images {\it C-A} and {\it C-B} at the top and middle,
respectively. Example light-echo positions are indicated with blue
(epoch {\it A}) and red (epoch {\it B} and {\it C}) arrows.  The
bottom right panel shows a zoom of the spectrograph slit, indicated
with a blue line.  For all panels north is up and east is to the left.
Applying the vector method that previously allowed us to identify the
source of the light echoes from the supernovae (SNe) that produced the
SN remnants 0509$-$67.5, Cas~A, and Tycho\cite{Rest05b,Rest08b}, we
find that a dramatic brightening of $\eta$~Car must be the origin.  In
these echoes, unlike those of Galactic SNe\cite{Rest11_leprofile},
there is still significant spatial overlap even at separations of one
light-year, suggesting that the duration of the event causing them
must be significantly longer than one year.  We also see brightening
of 2 magnitudes or more within 8 years.  Thus, the Lesser Eruption
from 1887 to 1896, which brightened by only a magnitude, is excluded
as the source.
}
{fig:diffim}

\fignature{lc_noblkline.pdf}{0.53}
{Historical and light-echo lightcurve of $\eta$~Car. The historical
light curve\cite{SF11} in visual apparent magnitudes is shown with
black circles and lines, with error bars indicating approximate
uncertainties in these eye estimates.  Light echo brightnesses (SDSS
$i$; error bars are the s.d.) from our eight modern images spanning
$\sim$8~years are displayed shifted by 174.2~yrs (green circles),
167.95~yrs (red circles), and 166.28~yrs (blue circles), in order to
illustrate the best-matching time delays for the 1838, 1843, and 1845
outbursts, respectively.
The first epoch is an upper limit indicated with an arrow. 
The upper panel shows the full time range of the Great
Eruption and therefore shows all three potential matches, whereas the
lower panels show the brightnesses from seven of our eight modern
epochs in a magnified time period around each peak.
}
{fig:lc}

\fignature{eta03_spec.pdf}{0.63}
{Light echo spectra of $\eta$~Car's Great Eruption.  The three optical
low-resolution spectra of the light echo (black lines) were taken at
J2000 position RA=10:44:12.127 and Decl.=$-$60:16:01.69 in March and
April 2011 obtained at the Magellan~I 6.5-m and du~Pont 2.5-m
telescopes of the Las Campanas Observatory, Chile. 
A log of the spectroscopic observations and
details of the spectra is presented in Table~S1.
\ifcheck
  \redpen{CHECK} 
\fi
The slit positions differ only slightly in slit angle.
The spectra were reduced using standard techniques and then
wavelength-calibrated using observations of an HeNeAr lamp. The
wavelength calibration was checked and corrected using night-sky
emission lines, especially \ion{O}{I} 5577\AA, and OH lines in the red
part of the spectrum. We flux-calibrated the spectra using a flux
standard observed the same night as the science observations.  The
left panel shows the spectra from 5000 to 9000\AA.  The spectra are
not corrected for reddening nor for the blue-ward scattering by the
dust.  The blue lines show for comparison spectra of three examples of
SN imposters: SN~1997bs, SN~2009ip, and UGC~2773-OT1.  The right panel
shows the H$\alpha$ and [\ion{N}{II}] emission lines. Note that the
background emission-line subtraction is incomplete since they
spatially vary. Also, for
\specMhi\ H$\alpha$ is at the edge of the chip and is therefore uncertain.
}
{fig:spec}

\fignature{hrdiagram2.pdf}{0.8}
{HR diagram with LBVs and $\eta$~Car.
Adaptation\cite{Smith04} of an HR diagram showing LBVs, related
hypergiant stars, and the peak luminosities of LBV-like eruptions. The
grey bands denote the typical locations of LBVs in quiescence and in S
Doradus excursions.  Temperatures for the Great Eruption and 1890
eruption of $\eta$~Car are based on the echo spectra presented here
and the F-type spectrum of the 1890 event\cite{Walborn77},
respectively.
%Values for
%UGC2773-OT and SN~2009ip are taken from Smith et al.\cite{Smith10}.
The temperature of 10,000~K for SN~2009ip is based on the observed
continuum shape, but this is only a lower limit because of the
possible effects of circumstellar or host galaxy
reddening\cite{Smith10}.  Because of the presence of \ion{He}{I} lines
in the spectrum, the true temperature is probably much hotter.
%The 7000~K temperature of UGC2773-OT indicated by its continuum
%temperature is a lower limit for the same reason, and the true
%temperature is also likely much hotter
The 8500~K temperature of UGC2773-OT is indicated by the F-type
absorption features in its spectrum, and this temperature is
relatively independent of reddening\cite{Smith10,Foley11}.
}
{fig:hr}
